% ===================================================================
%  TABEL Nr. 6 — Die huis van binne, die meubels, die voeding,
%  die verligting.
%  Meme regime que les precedents.
%
%  LE RENVOI (16) DU BLOC t06-c1-06-1 EST CELUI DU FRANCAIS, NON CELUI
%  DE L'IDO : le fac-simile ido y ouvre sur « (6) », qui est le savon.
%  C'est l'un des trois ecarts declares de outils/renvoji.py, et les
%  traductions le rendent, comme le neerlandais et le tcheque.
%
%  ON MODERNISE LA LANGUE, JAMAIS LES CHOSES. Le « vekigilo » est un
%  reveil a remonter, la « bujio » une kers, le « petrol-lampo » une
%  paraffienlamp. On n'ecrit pas un mot d'aujourd'hui pour un objet de
%  1926.
%
%  ET LA CUISINE EST LE PIEGE DE CETTE COLONNE : c'est le vocabulaire
%  ou l'afrikaans et le neerlandais partagent le plus de formes, et ou
%  la difference tient a une lettre — schotel/skottel,
%  schoorsteen/skoorsteen, zeep/seep, zout/sout, zeef/sif.
% ===================================================================

%%K t06-apar-1 apar
\VUsaut{6.0mm}
\VUtitre{15.0pt}{60}{TABEL Nr. 6}
\VUsaut{1.4mm}
\VUtitre{11.4pt}{0}{Die huis van binne, die meubels, die voeding,}
\VUsaut{0.4mm}
\VUtitre{11.4pt}{0}{die verligting.}
\VUsaut{2.2mm}

%%K t06-apar-2 apar
Die huis is klaar. Jan se oom het in sy nuwe woning ingetrek, waarvan ons die indeling
ken. Kom ons kyk nou hoe hy sy huis gemeubileer het. Ons sal eers verken:

%%K t06-tit-1 sub
\VUpk{10.2pt}{40}{Die slaapkamer.}
\VUsaut{3.0mm}

%%K t06-c1-01-1 p
1. --- Ons kom ongeleë, want die kamermeisie is besig om die \VUgras{kamer} skoon te
maak.

%%K t06-c1-02-1 p
2. --- Op die \VUgras{waskastjie} \textsuperscript{(1)} is hier en daar verskillende
voorwerpe waarmee 'n mens hom was, hom kam, kortom sy toilet maak. Dit is die
\VUgras{waskom} \textsuperscript{(2)} en die \VUgras{waterkan} \textsuperscript{(3)},
die \VUgras{haarborsel} \textsuperscript{(4)}, die \VUgras{kam} \textsuperscript{(5)},
die \VUgras{seep} \textsuperscript{(6)} en die \VUgras{spons} \textsuperscript{(7)},
laastens 'n \VUgras{flessie} \textsuperscript{(8)} vir die vloeibare tandmiddel.

%%K t06-c1-03-1 p
3. --- Op die vloer, voor die waskastjie, hier is die \VUgras{emmer}
\textsuperscript{(9)} om die vuil water op te vang.

%%K t06-c1-04-1 p
4. --- Ons sien in die \VUgras{voorkamer} \textsuperscript{(10)} 'n paar
\VUgras{klerehake} \textsuperscript{(13)} en twee \VUgras{sambrele}
\textsuperscript{(11)} in die \VUgras{sambreelstaander} \textsuperscript{(12)}.

%%K t06-c1-05-1 p
5. --- Aan die ander kant van die deur is die \VUgras{spieëlkas} \textsuperscript{(14)}
wat die \VUgras{linne} \textsuperscript{(15)} bevat.

%%K t06-c1-06-1 p
6. --- Kom ons maak gebruik daarvan dat die \VUgras{kamermeisie}
\textsuperscript{(16)} die \VUgras{bed} \textsuperscript{(17)} opmaak, om die
verskillende dele daarvan te verken. Die \VUgras{bedkatel} \textsuperscript{(20)} bevat
'n goeie \VUgras{veermatras} \textsuperscript{(21)} waarop Justina dadelik 'n wol
\VUgras{matras} \textsuperscript{(22)} sal sit. Daarna sal sy van die stoele die twee
\VUgras{lakens} \textsuperscript{(23)} en die \VUgras{kombers} \textsuperscript{(24)}
neem. Dan sal sy aan die koppenent die \VUgras{dwarskussing} \textsuperscript{(25)} en
die \VUgras{kopkussing} \textsuperscript{(26)} sit wat saam met die \VUgras{veerkombers}
\textsuperscript{(27)} op die \VUgras{bedmatjie} \textsuperscript{(32)} lê. Sy gebruik 'n
veerkwas om die \VUgras{gordyne} \textsuperscript{(19)} en die \VUgras{bedhemel}
\textsuperscript{(18)} af te stof en hier is 'n deel van die werk klaar.

%%K t06-c1-07-1 p
7. --- Dan sal sy die \VUgras{bedkassie} \textsuperscript{(28)} 'n bietjie moet
opruim, die \VUgras{wekker} \textsuperscript{(29)} moet opwen, die \VUgras{naglampie}
\textsuperscript{(30)} moet gereedmaak, die \VUgras{kers} \textsuperscript{(31)} moet
afhaal om die kandelaar skoon te maak.

%%K t06-tit-2 sub
\VUsaut{5.0mm}
\VUpk{10.2pt}{40}{Die werkmandjie en die kaggeltoebehore.}
\VUsaut{3.0mm}

%%K t06-c1-08-1 p
8. --- Die \VUgras{werkmandjie} \textsuperscript{(34)} naby die \VUgras{bankie}
\textsuperscript{(33)} het ook baie nodig om opgeruim te word. Sy sal die
\VUgras{borduurwerk} \textsuperscript{(35)} moet opvou, in die
\VUgras{werktafeltjie} \textsuperscript{(38)} die \VUgras{vingerhoed}, die
\VUgras{gare} \textsuperscript{(36)}, die \VUgras{naalde} \textsuperscript{(37)} en die
\VUgras{skêr} \textsuperscript{(39)} moet toemaak en die deksel van die tafeltjie met die
\VUgras{sleutel} \textsuperscript{(40)} moet sluit.

%%K t06-c1-09-1 p
9. --- Op die \VUgras{eenpoottafeltjie} \textsuperscript{(41)} in die middel sal Justina
die water van die \VUgras{karaf} \textsuperscript{(42)} vernuwe. Sy sal
\VUgras{suiker} in die \VUgras{suikerpot} \textsuperscript{(43)} sit, die
\VUgras{suikertangetjie} \textsuperscript{(44)} skoonmaak en die \VUgras{klereborsel}
\textsuperscript{(46)} wegneem.

%%K t06-c1-10-1 p
10. --- Die regterkant van die kamer sal minder werk van haar vra, want die \VUgras{lang
leunstoel} \textsuperscript{(47)} en sy \VUgras{kussing} \textsuperscript{(48)} is op
hulle regte plek, en omdat dit nie koud weer is nie, is niks verskuif onder die
toebehore van die \VUgras{kaggel} \textsuperscript{(49)} nie. Die \VUgras{skopgraaf}
\textsuperscript{(50)}, die \VUgras{tang} \textsuperscript{(51)}, die
\VUgras{vuurhonde} \textsuperscript{(55)}, die \VUgras{asrand} \textsuperscript{(56)} is
skoon; die \VUgras{vuurskerm} \textsuperscript{(54)} is nie verskuif nie en die
\VUgras{houtkis} \textsuperscript{(52)} is goed van \VUgras{blokke}
\textsuperscript{(53)} voorsien.

%%K t06-noto-1 noto
\VUnotes{143.2mm}{%
(*) Babo = klein kindjie, die E. baby, F. bébé, I. bambino.
\par
}

%%K t06-noto-2 noto
\VUnotes{143.2mm}{%
(*) Lambrequino = F. lambrequin, S. lambrequines, Port. lambrequins, selfs D. E.
lambrequins, dus D. E. F. P. S.%
}

%%K t06-c1-11-1 p
11. --- Tog sal sy die \VUgras{slingerhorlosie} \textsuperscript{(57)}, die
\VUgras{kandelabers} \textsuperscript{(58)} en die \VUgras{sakgoeddosies}
\textsuperscript{(59)} moet afstof, en laastens die \VUgras{spieël}
\textsuperscript{(60)} moet afvee.

%%K t06-c1-12-1 p
12. --- Wat die \VUgras{wieg} \textsuperscript{(61)} van die kindjie (van die babatjie
(*)) betref, dit is reeds gereed.

%%K t06-tit-3 sub
\VUsaut{5.0mm}
\VUpk{10.2pt}{40}{Die sitkamer.}
\VUsaut{3.0mm}

%%K t06-c2-01-1 p
1. --- Nou hier is die sitkamer. Dit is 'n baie mooi \VUgras{kamer}. Die kaggel, wat in
die regterhoek staan, is met 'n fyn \VUgras{beeldjie} \textsuperscript{(1)} versier, met
twee mooi \VUgras{vase} \textsuperscript{(3)} en met 'n fluweel \VUgras{lambrekyn}
\textsuperscript{(2)} (*) gedrapeer.

%%K t06-c2-02-1 p
2. --- Om te keer dat die vonke van die vuurherd die mat aan die brand steek, het 'n
mens 'n koper \VUgras{vonkskerm} \textsuperscript{(4)} voor die vuurherd gesit.

%%K t06-c2-03-1 p
3. --- Naby is 'n elegante \VUgras{skryftafeltjie} \textsuperscript{(5)} vir dames oop.
Daarop skryf die \VUgras{vrou van die huis} \textsuperscript{(23)} haar briewe.

%%K t06-c2-04-1 p
4. --- Die venster is met \VUgras{glasgordyne} \textsuperscript{(6)} met
\VUgras{koorde} \textsuperscript{(8)} versier wat aan die \VUgras{haakjies}
\textsuperscript{(7)} vasgehaak is, en met stof-oorgordyne wat \VUgras{bo gedrapeer}
\textsuperscript{(9)} is.

%%K t06-c2-05-1 p
5. --- In die vensternis bevat 'n \VUgras{potbedekker} \textsuperscript{(11)} 'n groen
\VUgras{plant} \textsuperscript{(10)}, 'n pragtige palm waarvan die blare hulle byna tot
by die \VUgras{paraffienlamp} \textsuperscript{(12)} uitstrek.

%%K t06-c2-06-1 p
6. --- Die \VUgras{lampskerm} \textsuperscript{(13)} van die lamp is deur Maria
gereëlgemaak, die bekoorlike \VUgras{jongmeisie} \textsuperscript{(14)} wat by die
klavier \textsuperscript{(15)} sit. Sy sit baie regop op die \VUgras{bankie}
\textsuperscript{(16)} en studeer 'n musiekstuk. Sy is baie bekwaam en sorgsaam, soos
haar \VUgras{musiekkas} \textsuperscript{(17)} bewys, wat volmaak opgeruim is.

%%K t06-tit-4 sub
\VUsaut{5.0mm}
\VUpk{10.2pt}{40}{Die deugniet.}
\VUsaut{3.0mm}

%%K t06-c2-07-1 p
7. --- Haar broer, Paulus, soek 'n \VUgras{boek} \textsuperscript{(19)} in die
\VUgras{boekrak} \textsuperscript{(18)}. Hy is 'n \VUgras{jongmens}
\textsuperscript{(20)}, beweeglik en ondraaglik. Terwyl hy hom aan die boekrak vashaak om
die boek te bereik wat te hoog is, waag hy dit om die \VUgras{borsbeeld}
\textsuperscript{(21)} wat bo staan te laat val.

%%K t06-c2-08-1 p
8. --- Hy maak gebruik daarvan, om daardie dwaasheid te doen, dat sy moeder besig is om
met 'n vriendin te gesels, 'n \VUgras{dame} \textsuperscript{(22)} wat 'n paar dae in
haar huis kom deurbring het. Die moeder sit op die \VUgras{rusbank}
\textsuperscript{(24)} naby die \VUgras{leunstoel} \textsuperscript{(25)}, en haar
vriendin is op die \VUgras{sitplek} \textsuperscript{(27)} waarvan 'n mens net die
\VUgras{rugleuning} \textsuperscript{(26)} sien.

%%K t06-c2-09-1 p
9. --- Die groot \VUgras{raam} \textsuperscript{(30)} wat teen die muur hang, bevat die
\VUgras{portret} \textsuperscript{(29)} van 'n familielid. In die \VUgras{album}
\textsuperscript{(32)} wat op die tafel gesit is, is die \VUgras{foto's}
\textsuperscript{(33)} van die ander gesinslede bymekaargebring. Die kinders het dit pas
deurgeblaai, want die \VUgras{stoele} \textsuperscript{(31)} staan nog om die tafel
saam.

%%K t06-c2-10-1 p
10. --- Die \VUgras{Chinese vaas} \textsuperscript{(34)} van porselein, wat naby die
\VUgras{papiermes} \textsuperscript{(35)} staan, is aan Maria vir haar naamdag geskenk,
en die \VUgras{skerm} \textsuperscript{(36)} wat die hoek van die kamer versier, is deur
haar oom uit China gebring.

%%K t06-c2-11-1 p
11. --- Vrolik gemaak deur die pragtige \VUgras{skilderye} \textsuperscript{(37)} wat
teen die \VUgras{muur} \textsuperscript{(42)} vasgehaak is, verlig deur die
\VUgras{plafonkroonlugter} \textsuperscript{(38)}, waarvan die \VUgras{elektriese
lampe} \textsuperscript{(39)} saans vrolik skyn, lyk daardie sitkamer baie aangenaam. Die
sagte \VUgras{mat} \textsuperscript{(40)} wat op die parket uitgesprei is, en die so
maklike \VUgras{elektriese klokkie} \textsuperscript{(41)} voltooi sy gerief.

%%K t06-tit-5 sub
\VUsaut{3.6mm}
\VUpk{10.2pt}{40}{Die eetkamer.}
\VUsaut{3.0mm}

%%K t06-c3-01-1 p
1. --- Die etensklok het gelui. Die hele gesin, behalwe Maria, is om die tafel
bymekaar. Die eetkamer is aangenaam verwarm deur die uitstekende \VUgras{stoof}
\textsuperscript{(1)} van erdewerk, met 'n dun \VUgras{pyp} \textsuperscript{(2)}.

%%K t06-c3-02-1 p
2. --- Daardie kamer het nie veel meubels nie. Langs die muur hang, bo die
\VUgras{bankie} \textsuperscript{(4)}, 'n sierlike \VUgras{fonteintjie}
\textsuperscript{(3)}. Baie naby is die \VUgras{buffet} \textsuperscript{(5)} waarop 'n
mens die koue geregte sit, die nagereg-skottels, en die verskillende tafelgerei wat 'n
mens nie dadelik nodig het nie.

%%K t06-c3-03-1 p
3. --- So sien ons op die boonste plank 'n \VUgras{gebraaide hoender}
\textsuperscript{(7)}, 'n \VUgras{vis} \textsuperscript{(8)} en 'n \VUgras{skaal}
\textsuperscript{(10)} met \VUgras{vrugte} \textsuperscript{(9)}. Onder is die
\VUgras{kaas} \textsuperscript{(11)}, die \VUgras{brood} \textsuperscript{(12)}, die
\VUgras{flessiestaander} \textsuperscript{(13)} wat alles bevat wat nodig is om die
slaai aan te maak: die olie, die asyn, die \VUgras{sout} \textsuperscript{(14)} en die
\VUgras{peper} \textsuperscript{(15)}. Laastens, op die onderste plank is 'n
\VUgras{koek} \textsuperscript{(17)} naby die \VUgras{mosterdpotjie}
\textsuperscript{(16)}.

%%K t06-c3-04-1 p
4. --- Die eetkamerhorlosie, wat 'n mens die \VUgras{muurhorlosie}
\textsuperscript{(20)} noem, het 'n baie eienaardige vorm. Dit is in 'n houtraam ingebou
wat boonop 'n \VUgras{barometer} \textsuperscript{(19)} en 'n \VUgras{termometer}
\textsuperscript{(18)} bevat.

%%K t06-tit-6 sub
\VUsaut{4.6mm}
\VUpk{10.2pt}{40}{Die bediening van die ete.}
\VUsaut{3.0mm}

%%K t06-c3-05-1 p
5. --- Teen al die mure van die kamer is \VUgras{houtpanele} \textsuperscript{(21)} van
gewaste eikehout aangebring. 'n Stof \VUgras{deurgordyn} \textsuperscript{(22)} sluit
die deur wat toegang tot die veranda gee, goed genoeg. Daar sal die gesin na die ete
koffie gaan drink. Die \VUgras{koppies} \textsuperscript{(24)} en die \VUgras{pierings}
\textsuperscript{(25)} is reeds gereed op die \VUgras{skottelkas}
\textsuperscript{(23)}.

%%K t06-c3-06-1 p
6. --- Bo die tafel is 'n \VUgras{hanglamp} \textsuperscript{(26)} aan die plafon
vasgemaak waarvan die baie helder lamp saans die hele eetkamer verlig.

%%K t06-c3-07-1 p
7. --- Kom ons verken nou die \VUgras{eettafel} \textsuperscript{(27)} self. Toe die
gesin ingekom het, was die tafelgerei gereed. Op die \VUgras{tafeldoek}
\textsuperscript{(28)} is, op die plek van elke tafelgenoot, 'n \VUgras{bord}
\textsuperscript{(33)}, 'n \VUgras{servet} \textsuperscript{(29)}, 'n \VUgras{lepel}
\textsuperscript{(34)}, 'n \VUgras{vurk} \textsuperscript{(35)}, 'n \VUgras{mes}
\textsuperscript{(36)} en 'n \VUgras{glas} \textsuperscript{(32)} gesit. Daar was ook
\VUgras{bottels} \textsuperscript{(30)} vir wyn en 'n \VUgras{karaf}
\textsuperscript{(31)} vir water.

%%K t06-c3-08-1 p
8. --- Die \VUgras{bediende} \textsuperscript{(39)} het eers die \VUgras{sopkom}
\textsuperscript{(37)} gebring, waarin nog 'n bietjie sop damp. Nou bedien hy die eerste
\VUgras{gereg} \textsuperscript{(38)}, 'n vleisgereg. Daarna sal hy dié aanbied wat ons
reeds genoem het, laastens, na die \VUgras{nagereg}, sal hy gebruik maak daarvan dat sy
werkgewers na die veranda gegaan het, om die tafelgerei weg te neem en die eetkamer weer
op te ruim.

%%K t06-c3-08-2 p
\textit{Nota.} --- \textit{Die gewone woorde wat op die voeding betrekking het, word hier
baie beknop aangedui. Die lys sal op die twee tabelle « Die hotel » (nr. 13) en « Die
mark » (nr. 15) aangevul word.}

%%K t06-tit-7 sub
\VUsaut{4.6mm}
\VUpk{10.2pt}{40}{Die kombuis.}
\VUsaut{3.0mm}

%%K t06-c4-01-1 p
1. --- In die kombuis, waarna ons deur die half oop deur begin kyk, sien ons 'n
skilderagtige wanorde. Voor die \VUgras{vuurherd} \textsuperscript{(1)}, waarin 'n
\VUgras{vuur} \textsuperscript{(3)} knetter, is 'n \VUgras{braaispit}
\textsuperscript{(5)} opgestel. 'n Mooi \VUgras{braaistuk} \textsuperscript{(7)} is
\VUgras{aan die spit gesteek} \textsuperscript{(6)} en braai gaar.

%%K t06-c4-02-1 p
2. --- Die \VUgras{blaaspyp} \textsuperscript{(8)} en die \VUgras{koolskopgraaf}
\textsuperscript{(9)}, wat 'n mens pas gebruik het, het op die vloer bly lê net soos die
\VUgras{blokke} \textsuperscript{(10)}.

%%K t06-c4-03-1 p
3. --- Die bokant van die kaggel is opgeruim; die \VUgras{lantern}
\textsuperscript{(11)}, die \VUgras{vuurhoutjiedosie} \textsuperscript{(13)}, waarin die
\VUgras{vuurhoutjies} \textsuperscript{(12)} is, die \VUgras{kandelaar}
\textsuperscript{(14)} en die \VUgras{blaker} \textsuperscript{(15)} met 'n
\VUgras{kers} \textsuperscript{(16)} is op hulle plek. Maar die groot \VUgras{koolemmer}
\textsuperscript{(18)}, vol \VUgras{steenkool} \textsuperscript{(17)} wat die
\VUgras{kokkin} \textsuperscript{(23)} pas in die kelder volgemaak het, het in die middel
van die kombuis bly staan.

%%K t06-c4-04-1 p
4. --- Waarlik, die arme vrou moet haas sodat die ete betyds gereed is. Nou is sy by die
\VUgras{kookstoof} \textsuperscript{(19)} en roer 'n \VUgras{sous}
\textsuperscript{(24)} wat nie te veel moet aanbrand nie.

%%K t06-c4-05-1 p
5. --- In die \VUgras{oond} \textsuperscript{(20)} bak 'n koek wat sy vir die kinders se
tussenete voorberei het. Bo die kookstoof hang die \VUgras{potlepel}
\textsuperscript{(21)} en die \VUgras{skuimspaan} \textsuperscript{(22)}.

%%K t06-tit-8 sub
\VUsaut{4.0mm}
\VUpk{10.2pt}{40}{Die vaatwerk.}
\VUsaut{3.0mm}

%%K t06-c4-06-1 p
6. --- Die \VUgras{kookgerei} \textsuperscript{(25)}, wat agterin die kamer opgehang is,
is baie skoon. Die mooi koper \VUgras{kastrolle} \textsuperscript{(26)}, die
\VUgras{melkkan} \textsuperscript{(27)}, die \VUgras{sif} \textsuperscript{(28)} en die
\VUgras{rasper} \textsuperscript{(29)} is sorgvuldig skoongemaak.

%%K t06-c4-07-1 p
7. --- Die \VUgras{soutbakkie} \textsuperscript{(30)} is op die skottelkas gesit naby
die \VUgras{teepot} \textsuperscript{(31)} en die \VUgras{oliefles}
\textsuperscript{(32)}. Die \VUgras{laai} \textsuperscript{(33)} bevat waarskynlik die
eetgerei en die kombuismesse.

%%K t06-c4-08-1 p
8. --- Die \VUgras{besem} \textsuperscript{(34)} en die \VUgras{veerkwas}
\textsuperscript{(35)} staan in 'n hoek. Die kokkin het pas die \VUgras{kapblok}
\textsuperscript{(36)} gebruik, sy het dit naby die venster laat staan.

%%K t06-c4-09-1 p
9. --- 'n \VUgras{Handdoek} \textsuperscript{(37)} hang teen die muur, naby die
\VUgras{tregter} \textsuperscript{(38)}. In daardie deel van die kombuis is die
\VUgras{rooster} \textsuperscript{(39)}, die \VUgras{kapmes} \textsuperscript{(40)} en
die \VUgras{kapplank} \textsuperscript{(41)} weggesit. Verder, bo die kapmes, op 'n
\VUgras{plank} \textsuperscript{(42)}, is daar \VUgras{potte} \textsuperscript{(43)},
die \VUgras{kookpot} \textsuperscript{(44)} met sy \VUgras{deksel}
\textsuperscript{(45)} en die \VUgras{ketel} \textsuperscript{(46)}. Die \VUgras{pan}
\textsuperscript{(47)} hang naby.

%%K t06-c4-10-1 p
10. --- Net die koper \VUgras{kraan} \textsuperscript{(48)} werp 'n glans in hierdie
hoek. Die \VUgras{opwasbak} \textsuperscript{(49)} waarop die dik \VUgras{kruik}
\textsuperscript{(50)} gesit is, is waarskynlik baie gerieflik met sy klein kassie,
waarin 'n mens die \VUgras{asynvaatjie} \textsuperscript{(51)} met sy \VUgras{kraantjie}
\textsuperscript{(52)} bewaar, die \VUgras{afvalbak} \textsuperscript{(53)} en die
\VUgras{vuil borde} \textsuperscript{(54)}.

%%K t06-tit-9 sub
\VUsaut{4.0mm}
\VUpk{10.2pt}{40}{Ander voorwerpe in die kombuis.}
\VUsaut{3.0mm}

%%K t06-c4-11-1 p
11. --- Die \VUgras{balie} \textsuperscript{(55)} waarin 'n mens die \VUgras{groente}
\textsuperscript{(58)} was, en die gietyster \VUgras{ketel} \textsuperscript{(56)} wat op
die \VUgras{teëlvloer} \textsuperscript{(57)} gesit is, het waarskynlik ook hulle plek in
daardie kassie.

%%K t06-c4-12-1 p
12. --- Die kokkin kom sekerlik nie stowe kort nie, want hier is boonop, op die hoek van
die tafel, 'n \VUgras{gasstofie} \textsuperscript{(59)} waarop water in 'n klein kastrol
warm word. Die \VUgras{stoom} \textsuperscript{(60)} wat daaruit ontsnap, dui vir ons aan
dat dit waarskynlik kook. Vermoedelik sal hierdie water vir die koffie gebruik word, want
hier is op die ander punt van die tafel die \VUgras{koffiekan} \textsuperscript{(64)} en
die \VUgras{koffiemeul} \textsuperscript{(63)}.

%%K t06-c4-13-1 p
13. --- Die stofie is met die \VUgras{gasbrander} \textsuperscript{(62)} verbind deur 'n
lang \VUgras{rubberpypie} \textsuperscript{(61)}.

%%K t06-c4-14-1 p
14. --- Die \VUgras{dagmeisie} \textsuperscript{(66)} wat die kokkin kom help, maak die
groente vir die ete gereed. Wanneer hulle skoongemaak en gewas is, sal sy hulle in die
\VUgras{bak} \textsuperscript{(65)} sit terwyl sy op die uur wag om hulle te kook.

%%K t06-tit-10 sub
\VUsaut{4.0mm}
{\centering\fontsize{10.2pt}{10.2pt}\selectfont\textls[40]{\scshape Die badkamer.}
 \textit{(Die baaikamer.)}\par}
\VUsaut{3.0mm}

%%K t06-c5-01-1 p
1. --- Ons het nog net een onbeskeidenheid oor om die vernaamste kamers van die woning
te leer ken: om die inrigting van die badkamer te sien. Dit sal nie lank duur nie, want
dit is nie groot nie, en ons sal gou weet dat die \VUgras{bad} \textsuperscript{(1)} 'n
goeie \VUgras{badverwarmer} \textsuperscript{(2)} het, dat die \VUgras{seepbakkie}
\textsuperscript{(3)} binne bereik van die badende se hand is en dat die
\VUgras{handdoekrak} \textsuperscript{(5)} die \VUgras{handdoeke} \textsuperscript{(4)}
sekerlik maklik sal laat droog word.

%%K t06-c5-02-1 p
2. --- Daardie kamer het sy mure met \VUgras{teëls} \textsuperscript{(6)} bedek, wat dit
moontlik gemaak het om 'n \VUgras{stortapparaat} \textsuperscript{(7)} te installeer
sonder om enigsins te vrees dat die water, wat tydens sy gebruik terugspat, die mure sal
beskadig. 'n Sink \VUgras{bakkie} \textsuperscript{(8)} wat vir die voetbaddens dien,
voltooi die meubels.
\VUsaut{6.0mm}
\VUfilet{22mm}
